\chapter{Concept Generation}
\todo{kind of the explanation of our DOT and how we arrived at the concepts below, then we can specifically desribe the systems in the system level design options}


\section{Design Space Overview}

The design space was structured through three main Design Option
Trees (DOTs): aerial platform, detection and tracking, and balloon
interaction strategy. These three areas are strongly coupled and
together define the concepts carried into the conceptual design
phase. The full DOT is provided in the Baseline Report
(Appendix~A.4).

\subsection*{DOT 1 — Aerial Platform}
The aerial platform determines payload capacity,
controllability near the target, range, altitude capability,
energy demand, acoustic performance, and reusability.
\textbf{Rotary-wing} platforms are identified as the strongest
baseline option due to their hover capability, low-speed
controllability, and precedent in counter-UAS and aerial
manipulation systems. \textbf{Hybrid VTOL} configurations
(tilt-rotor, tail-sitter) are retained because of their
improved range and endurance, at the cost of increased
mechanical complexity. \textbf{Conventional fixed-wing}
platforms are retained for their endurance advantage.
Tilt-wing, rocket, flapping-wing, and coaxial helicopter
configurations were pruned for reasons of cost, complexity,
or certification risk.

\subsection*{DOT 2 — Detection and Tracking}
The detection and tracking architecture must acquire a
slow-moving, low radar-cross-section target and hand it to a
terminal tracker with sufficient accuracy for close-range
interaction. The most promising architecture is a
\textbf{hybrid-cued} approach: coarse target information from
a ground observer or external monitoring source is handed over
to onboard optical and/or thermal sensing for terminal
interception. RGB cameras (YOLO-family algorithms) are retained
as the primary onboard sensor for target confirmation and
tracking in favourable conditions. LWIR thermal imaging is
retained as a complementary sensor for low-light operation.
LiDAR is retained as an optional terminal sensor only.
Onboard radar\todo{double check}, IRST, and sonar were pruned on grounds of cost,
weight, or insufficient performance against a latex balloon
envelope.

\subsection*{DOT 3 — Balloon Interaction Strategy}
The interaction DOT is structured around three sequential
decision nodes: (1)~first contact and handling,
(2)~method to initiate descent, and (3)~method to manage
descent.

\paragraph{Node 1 — First contact.}
Net-based approaches (cooperative multi-UAV net, radially
deployed net, hanging-net sweep) and tether/lasso approaches
(line-snag V-yoke, cooperative tether closure) are retained
as the primary non-direct contact options. Direct contact
options retained include enveloping soft grippers, rigid
clamps for tether engagement, gecko-pad arrays, regular
adhesive pads, and active pneumatic suction. Options pruned
include talon/pincer claws (uncontrolled deflation risk),
articulated robotic arms (cost/complexity), cooperative
clamping (bladestrike risk), passive suction cups (insufficient grip force), and electromagnetic approaches (no guaranteed ferromagnetic target material).

\paragraph{Node 2 — Descent initiation.}
Retained options include integrated thrust (bi-directional
propeller or throttle-below-hover), ground-connected tether, controlled venting via an active venting module or insert valve. Uncontrolled deflation methods (directed energy, projectile, propeller puncture) were screened out for violating the non-pyrotechnic and debris constraints. Solid ballast was pruned for uncontrolled descent risk.

\paragraph{Node 3 — Descent management.}
No options are pruned at this node as all are dependent on
the selected platform and interaction concept. Options
retained include parachutes, drag surfaces, active drag
orientation, increased rotor thrust (UAS tow), lifting
surfaces (parafoil), gas envelope, and variable buoyancy
systems.


\section{Overview of the five selected concepts}


Table~\ref{tab:strawman_concept_overview} gives a compact overview of the candidate concepts selected from the retained design-option-tree branches. Each concept is described as a complete mission architecture combining an aerial platform, interaction method, and descent or recovery strategy. The descriptions provide a consistent basis for comparing the concepts in the trade-off. 

\begin{table}[H]
    \centering
    \small
    \caption{Candidate straw-man concepts entering the trade-off.}
    \label{tab:strawman_concept_overview}
    \begin{tabularx}{\textwidth}{@{}p{0.04\textwidth} p{0.18\textwidth} X X X@{}}
        \toprule
        \textbf{ID} 
        & \textbf{Concept name} 
        & \textbf{Aerial platform} 
        & \textbf{Interaction principle} 
        & \textbf{Descent/recovery principle} \\
        \midrule

        A 
        & Bottom Fishing
        & Four-propeller tail-sitter VTOL UAV with cruise and hover capability.
        & The UAV hovers below the payload and launches an upward capture net that closes around the payload through a tethered pouch mechanism.
        & The UAV tows the captured balloon-payload system downward using thrust and vehicle weight, with emergency breakaway and parachute backup. \\

        B 
        & Pick and drop
        & Carrier eVTOL deploys a team of net-carrier multirotors and one helper drone near the target.
        & The net-carrier drones envelop the balloon-payload system using a two-stage cinching net.
        & The secured bundle separates from the drones and descends under a guided parafoil. \\

        C 
        & Mechanical Clamping
        & Lift-and-cruise eVTOL UAV station-keeping beneath the payload.
        & Deployable aluminium beams extend and pivot upward to mechanically secure the payload or its attachment.
        & The coupled system descends under UAV weight, with directed thrust used for descent control, stabilization, and landing guidance. \\

        D 
        & Cooperative Variable-Weight Payload Net
        & Three heavy-lift multirotors carrying a shared payload-capture net on long tethers.
        & The UAV team positions a net below the target and closes it around the payload.
        & Ballast and net hardware provide downward force, while the UAV formation stabilises and steers the captured system during descent. \\

        E 
        & Kinetic Bolas
        & Lift-and-cruise hybrid UAV with hover capability for close-range engagement.
        & A pneumatic launcher fires two weighted projectiles connected by a leader line, wrapping around the suspension line.
        & The UAV reels out a tow line and uses vehicle weight and thrust to guide the balloon-payload system toward recovery. \\

        \bottomrule
    \end{tabularx}
\end{table}
